% proposta.tex — Proposta Inicial · ForensiQ · UC 21184
% Compilar: latexmk -pdf -output-directory=../docs/report proposta.tex
\documentclass[a4paper,11pt]{article}

% ── Codificação e língua ──────────────────────────────────────────────
\usepackage[utf8]{inputenc}
\usepackage[T1]{fontenc}
% NOTA: descomentar quando texlive-lang-portuguese estiver instalado
% \usepackage[portuguese]{babel}

% ── Layout ────────────────────────────────────────────────────────────
\usepackage[margin=2.5cm]{geometry}
\usepackage{parskip}            % espaço entre parágrafos, sem indentação
\usepackage{setspace}
\onehalfspacing

% ── Tabelas e listas ─────────────────────────────────────────────────
\usepackage{booktabs}
\usepackage{longtable}
\usepackage{enumitem}
\setlist{nosep,leftmargin=1.5em}

% ── Fontes e tipografia ──────────────────────────────────────────────
\usepackage{lmodern}
\usepackage{microtype}

% ── Imagens ──────────────────────────────────────────────────────────
\usepackage{graphicx}

% ── Hyperlinks ────────────────────────────────────────────────────────
\usepackage[colorlinks=true,linkcolor=black,urlcolor=blue,citecolor=black]{hyperref}

% ── Cabeçalhos ────────────────────────────────────────────────────────
\usepackage{fancyhdr}
\pagestyle{fancy}
\fancyhf{}
\fancyhead[L]{\small ForensiQ — Proposta Inicial}
\fancyhead[R]{\small UC 21184 · 2025/26}
\fancyfoot[C]{\thepage}
\renewcommand{\headrulewidth}{0.4pt}

% ── Tradução de títulos automáticos (sem babel) ──────────────────────
\renewcommand{\contentsname}{Índice}
\renewcommand{\refname}{Referências}
\renewcommand{\tablename}{Tabela}
\renewcommand{\figurename}{Figura}

% ── Bibliografia ──────────────────────────────────────────────────────
\usepackage[round,authoryear]{natbib}
\bibliographystyle{plainnat}

% ══════════════════════════════════════════════════════════════════════
\begin{document}

% ── CAPA ──────────────────────────────────────────────────────────────
\begin{titlepage}
\centering
\vspace*{2cm}

{\Large\bfseries Universidade Aberta}\\[0.3cm]
{\large Departamento de Ciências e Tecnologia}\\[0.3cm]
{\large Licenciatura em Engenharia Informática}\\[2cm]

{\LARGE\bfseries ForensiQ}\\[0.4cm]
{\large Plataforma Modular de Gestão de Prova Digital\\para First Responders}\\[2.5cm]

{\large\bfseries Proposta Inicial}\\[0.3cm]
{\large UC 21184 — Projeto de Engenharia Informática}\\[1.5cm]

\begin{tabular}{rl}
\textbf{Estudante:} & João Rodrigues — 2203474 \\
\textbf{Orientador:} & Professor Pedro Duarte Pestana \\
\textbf{Data:} & 22 de março de 2026 \\
\textbf{Versão:} & 1.0 \\
\end{tabular}

\vfill
{\small 2.º Semestre · 2025/2026}
\end{titlepage}

% ── ÍNDICE (opcional, não conta para as 4 páginas) ───────────────────
\tableofcontents
\newpage

% ══════════════════════════════════════════════════════════════════════
% SINOPSE
% ══════════════════════════════════════════════════════════════════════
\section{Sinopse}

Os agentes que chegam primeiro a uma cena de crime — os \emph{first responders} —
recorrem actualmente a métodos manuais (papel, fotografias avulsas, anotações)
para registar e preservar a prova digital apreendida. Esta ausência de
padronização traduz-se em vulnerabilidades processuais reais: cadeia de custódia
inconsistente, prova sem localização verificável e procedimentos que variam de
agente para agente. O problema é identificado da experiência directa do estudante
como perito digital forense na Polícia de Segurança Pública (PSP).

O ForensiQ é uma aplicação web \emph{mobile-first} que digitaliza e padroniza
este processo, do momento da apreensão no terreno até à chegada ao laboratório e
remessa das conclusões da perícia. Distingue-se dos sistemas existentes (SAFE,
Axon Evidence) por ser desenhado para o agente sem formação técnica avançada, por
seguir os procedimentos da PSP, e por garantir integridade dos registos via
\emph{hash} SHA-256 conforme a norma ISO/IEC~27037
\citep{iso27037}. Diferencia-se do Projecto~\#38 da orientação
\citep{pestana2025} por não adquirir prova — apenas a gerir e controlar após
apreensão.

O resultado esperado é um protótipo funcional demonstrável num \emph{smartphone}
no terreno, com API REST documentada, módulo de forense digital completo e cadeia
de custódia com registo imutável verificável. O sucesso é medido pelos critérios
de aceitação observáveis definidos na Secção~\ref{sec:mvp}.

% ══════════════════════════════════════════════════════════════════════
% MVP
% ══════════════════════════════════════════════════════════════════════
\section{MVP — Funcionalidades e Critérios de Aceitação}
\label{sec:mvp}

O produto mínimo viável é composto por seis funcionalidades, cada uma com
critérios de aceitação observáveis e verificáveis em demonstração.

\subsection{F1 — Autenticação com perfis}

Dado email e password válidos de agente, o sistema autentica e devolve token JWT
em menos de 2~segundos; dado email inválido, devolve erro~401 sem expor
informação interna. Os perfis agente e perito forense digital têm permissões
distintas verificáveis.

\subsection{F2 — Registo de prova com fotografia e metadados GPS}

Um agente autenticado cria um registo de prova com fotografia, tipo, descrição e
localização GPS capturada automaticamente pelo \emph{browser}. O registo recebe
identificador único, \emph{timestamp} do servidor e \emph{hash} SHA-256 dos
metadados. Verificável num \emph{smartphone} em menos de 30~segundos.

\subsection{F3 — Cadeia de custódia com máquina de estados e log imutável}

A prova percorre estados em ordem fixa sem possibilidade de retrocesso:

\begin{center}
\small
\texttt{APREENDIDA $\rightarrow$ EM TRANSPORTE $\rightarrow$ RECEBIDA NO LABORATÓRIO $\rightarrow$ EM PERÍCIA $\rightarrow$ CONCLUÍDA $\rightarrow$ DEVOLVIDA/DESTRUÍDA}
\end{center}

Cada transição regista: agente responsável, \emph{timestamp} e \emph{hash} do
registo nesse momento. Nenhuma transição pode ser saltada. O registo é
\emph{append-only} — sem \texttt{UPDATE} nem \texttt{DELETE} na tabela de
custódia. Qualquer alteração directa na base de dados é detectável via
\emph{hash} \citep{casey2011,acpo2012}.

\subsection{F4 — Módulo de forense digital (ficha de dispositivo)}

O perito preenche ficha completa de dispositivo digital (tipo, marca, modelo,
estado, número de série) associada a uma prova existente.

\subsection{F5 — API REST documentada}

Mínimo 10 \emph{endpoints} funcionais documentados e testáveis via Swagger~UI
(drf-spectacular).

\subsection{F6 — Exportação de relatório de ocorrência em PDF}

Relatório gerado automaticamente com todos os registos de prova e histórico de
custódia da ocorrência. Acessível aos perfis agente e perito, permitindo ao
agente juntar o relatório ao expediente físico.

% ══════════════════════════════════════════════════════════════════════
% STACK TECNOLÓGICA
% ══════════════════════════════════════════════════════════════════════
\section{Stack Tecnológica}

\begin{table}[h]
\centering
\small
\begin{tabular}{@{}lll@{}}
\toprule
\textbf{Componente} & \textbf{Tecnologia} & \textbf{Justificação} \\
\midrule
Backend         & Django + DRF          & Estrutura convencional; autenticação e ORM built-in \\
Frontend        & HTML/CSS/JS vanilla   & Mobile-first sem overhead de framework; suficiente para MVP \\
Base de dados   & PostgreSQL            & Integridade referencial; append-only para logs de custódia \\
Autenticação    & JWT + RBAC            & Dois perfis: agente e perito forense digital \\
API docs        & drf-spectacular       & Documentação Swagger UI automática \\
Geolocalização  & Geolocation API + Leaflet.js & Sem API paga; funciona em telemóvel \\
Integridade     & SHA-256 (hashlib)     & Hash dos metadados conforme ISO/IEC~27037 \\
CI/CD           & GitHub Actions        & Pipeline de testes automáticos a cada push \\
\bottomrule
\end{tabular}
\caption{Stack tecnológica do ForensiQ.}
\end{table}

% ══════════════════════════════════════════════════════════════════════
% ARQUITECTURA C4 NÍVEL 1
% ══════════════════════════════════════════════════════════════════════
\section{Esboço de Arquitectura — C4 Nível 1 (Contexto)}

A Figura~\ref{fig:c4context} apresenta o sistema ForensiQ, os seus utilizadores
e o sistema externo com que interage, segundo a notação C4 \citep{brown2014}.

O \textbf{Agente} (first responder) regista provas no terreno via smartphone —
captura fotografia, localização GPS e metadados — e exporta relatórios em PDF
para juntar ao expediente físico. O \textbf{Perito forense digital} gere a
cadeia de custódia no laboratório, preenche fichas de dispositivo e exporta
relatórios. O único sistema externo é a \textbf{Browser Geolocation API}, que
fornece coordenadas GPS do dispositivo do agente.

\begin{figure}[h]
\centering
\includegraphics[width=0.85\textwidth]{figures/c4-context.png}
\caption{Diagrama de contexto C4 (Nível 1) do ForensiQ.}
\label{fig:c4context}
\end{figure}

O diagrama C4 de nível 2 (contentores: aplicação web, API REST e base de dados)
será produzido durante as semanas~3--4 e depositado em
\texttt{docs/architecture/}.

% ══════════════════════════════════════════════════════════════════════
% CALENDÁRIO
% ══════════════════════════════════════════════════════════════════════
\newpage
\section{Calendário Detalhado}

\begin{longtable}{@{}llp{7.5cm}l@{}}
\toprule
\textbf{Sem.} & \textbf{Datas} & \textbf{Conteúdo planeado} & \textbf{Marco} \\
\midrule
\endfirsthead
\toprule
\textbf{Sem.} & \textbf{Datas} & \textbf{Conteúdo planeado} & \textbf{Marco} \\
\midrule
\endhead

1 & 17--21 mar &
  Configuração do repositório a partir do template.
  Definição de stack e MVP. Envio de sinopse, MVP e calendário ao orientador. &  \\
\midrule

2 & 24--28 mar &
  Finalizar proposta inicial. Submeter na plataforma da UC.
  Aguardar feedback e aceitação do orientador. &
  \textbf{Proposta (25 mar)} \\
\midrule

3 & 31 mar--4 abr &
  Requisitos MoSCoW finais. Diagrama C4 nível 1 e nível 2.
  Modelo de dados ER (PostgreSQL). &  \\
\midrule

4 & 7--11 abr &
  ADRs obrigatórios (Django, PostgreSQL, SHA-256, máquina de estados, vanilla JS).
  Início da estrutura Django (\texttt{manage.py}, modelos). &  \\
\midrule

5 & 14--17 abr &
  Wireframes mobile-first (ecrãs principais).
  Implementação F1: autenticação JWT com perfis agente e perito. &  \\
\midrule

6 & 22--25 abr &
  Implementação F2: registo de prova com fotografia e GPS.
  Implementação F3: cadeia de custódia (máquina de estados, log append-only). &  \\
\midrule

7 & 28 abr--2 mai &
  Testes de integração dos módulos F1--F3. Preparação e execução da demo interna
  ao orientador (via Teams). &
  \textbf{Demo interna} \\
\midrule

8 & 5--6 mai &
  Redacção e submissão do relatório intercalar.
  Incorporação de feedback do orientador. &
  \textbf{Intercalar (6 mai)} \\
\midrule

9 & 7--9 mai &
  Implementação F4: módulo de forense digital (ficha de dispositivo).
  Início F5: endpoints da API REST. &  \\
\midrule

10 & 12--16 mai &
  Conclusão F5: API REST com 10+ endpoints + Swagger UI.
  Implementação F6: exportação de relatório PDF. &  \\
\midrule

11 & 19--23 mai &
  Testes de integração (Postman/Newman).
  Configuração GitHub Actions CI. HTTPS. &  \\
\midrule

12 & 26--30 mai &
  Testes end-to-end. Correcção de defeitos.
  Testes em dispositivo móvel real. &  \\
\midrule

13 & 2--6 jun &
  Polimento da interface. Início da redacção dos capítulos 4 e 5
  do relatório final. &  \\
\midrule

14 & 9--13 jun &
  Revisão completa do relatório final.
  Correcções de última hora no código e documentação. &  \\
\midrule

15 & 16--20 jun &
  Preparação da defesa oral (15 minutos). Ensaio com simulação de perguntas
  do júri. Reunião de preparação com orientador (Teams). &
  \textbf{Prep. defesa} \\
\midrule

16 & 24 jun &
  Submissão do relatório final e código na plataforma da UC. &
  \textbf{Final (24 jun)} \\

\bottomrule
\end{longtable}

% ══════════════════════════════════════════════════════════════════════
% REFERÊNCIAS
% ══════════════════════════════════════════════════════════════════════
\newpage
\bibliography{bibliography}

\end{document}
